\documentclass[12pt,a4paper]{article}
\usepackage[utf8]{inputenc}
\usepackage[margin=1in]{geometry}
\usepackage{titlesec}
\usepackage{tabularx}
\usepackage{hyperref}
\usepackage{enumitem}
\usepackage{graphicx}
\usepackage{listings}
\usepackage{color}
\usepackage{float} % Forces images to stay in place

% Code highlighting configuration
\definecolor{codegray}{rgb}{0.5,0.5,0.5}
\definecolor{codepurple}{rgb}{0.58,0,0.82}
\definecolor{backcolour}{rgb}{0.95,0.95,0.92}

\linespread{1.3} % Increase general line spacing for a fuller look

\lstset{
    backgroundcolor=\color{backcolour},   
    commentstyle=\color{codegray},
    keywordstyle=\color{blue},
    numberstyle=\tiny\color{codegray},
    stringstyle=\color{codepurple},
    basicstyle=\ttfamily\footnotesize,
    breakatwhitespace=false,         
    breaklines=true,                 
    captionpos=b,                    
    keepspaces=true,                 
    numbers=left,                    
    numbersep=5pt,                  
    showspaces=false,                
    showstringspaces=false,
    showtabs=false,                  
    tabsize=2,
    language=SQL
}

% Title Formatting
\titleformat{\section}{\normalfont\huge\bfseries}{}{0pt}{}
\titleformat{\subsection}{\normalfont\Large\bfseries}{}{0pt}{}

\begin{document}

% --- COVER PAGE ---
\begin{center}
    \vspace*{1cm}
    \textbf{\huge SVKM’s NMIMS} \\
    \vspace{0.5cm}
    \textbf{\Large School of Technology Management \& Engineering, Indore} \\
    \vfill
    \textbf{\large A.Y. 2025 - 26} \\
    \vspace{1cm}
    \textbf{\Large Course: Database Management Systems} \\
    \vspace{0.5cm}
    \textbf{\huge Project Report} \\
    \vfill
    
    \renewcommand{\arraystretch}{1.8}
    \begin{tabularx}{\textwidth}{|l|X|}
        \hline
        \textbf{Program} & B.Tech CE \\ \hline
        \textbf{Semester} & 4th \\ \hline
        \textbf{Name of the Project:} & \textbf{GreenChain: A Relational Database Framework for Quantifying Supply Chain Carbon Emissions} \\ \hline
        \textbf{Details of Project Members} & \\ \hline
    \end{tabularx}
    
    \vspace{1cm}
    
    \begin{tabularx}{\textwidth}{|l|l|X|}
        \hline
        \textbf{Batch} & \textbf{Roll No.} & \textbf{Name} \\ \hline
        A1 & F044 & Anupam Kanoongo \\ \hline
        A1 & F005 & Arjun Borkar \\ \hline
    \end{tabularx}
    
    \vfill
    \textbf{\Large Date of Submission:} \today
    \vspace*{1cm}
\end{center}

\newpage

% --- CONTRIBUTION PAGE ---
\section*{Project Member Contributions}

\renewcommand{\arraystretch}{2}
\begin{tabularx}{\textwidth}{|l|l|X|}
    \hline
    \textbf{Roll No.} & \textbf{Name:} & \textbf{Contribution} \\ \hline
    F044 & Anupam Kanoongo & Backend Architecture, Database Design, SQL Query Optimization, and System Integration. \\ \hline
    F005 & Arjun Borkar & Frontend Development (AngularJS), Data Visualization (Chart.js), and UI/UX Design. \\ \hline
\end{tabularx}

\vfill
\noindent \textbf{\Large Github link of your project:} \\
\url{https://github.com/AnupamKanoongo/DBMS_Project_GreenChain}

\vspace{1.5cm}
\textbf{Note:}
\begin{itemize}[itemsep=10pt]
    \item Create a readme file if you have multiple files
    \item All files must be properly named (Example: R004\_DBMSProject)
    \item Submit all relevant files of your work (Report, all SQL files, Any other files)
    \item Plagiarism is highly discouraged (Your report will be checked for plagiarism)
\end{itemize}

\vspace{1cm}
\textbf{Rubrics for the Project evaluation:}
\begin{table}[h]
    \centering
    \renewcommand{\arraystretch}{1.5}
    \begin{tabular}{|p{4cm}|p{4cm}|p{4cm}|c|}
        \hline
        First phase of evaluation & Innovative Ideas (5 Marks) & Design and Partial implementation (5 Marks) & 10 marks \\ \hline
        Final phase of evaluation & \multicolumn{2}{p{8cm}|}{Implementation, presentation and viva, Self-Learning and Learning Beyond classroom} & 10 marks \\ \hline
    \end{tabular}
\end{table}

\newpage

% --- INTERNAL TITLE PAGE ---
\begin{center}
    \vspace*{2cm}
    \section*{Project Report}
    \vfill
    \textbf{\huge GreenChain: A Relational Database Framework for Quantifying Supply Chain Carbon Emissions} \\
    \vfill
    \textbf{\Large Submitted by:} \\
    \vspace{1cm}
    \textbf{\Large Anupam Kanoongo, Roll number: F044} \\
    \textbf{\Large Arjun Borkar, Roll number: F005} \\
    \vfill
    \textbf{\Large Course: DBMS} \\
    \textbf{\Large AY: 2025-26}
    \vspace*{2cm}
\end{center}

\newpage

% --- TABLE OF CONTENTS ---
\section*{Table of Contents}
\vspace{1cm}
\renewcommand{\arraystretch}{2.2}
\begin{tabularx}{\textwidth}{|c|X|c|}
    \hline
    \textbf{Sr no.} & \textbf{Topic} & \textbf{Page no.} \\ \hline
    1 & Storyline & 4 \\ \hline
    2 & Components of Database Design & 5 \\ \hline
    3 & Entity Relationship Diagram & 6 \\ \hline
    4 & Relational Model & 7 \\ \hline
    5 & Normalization & 8 \\ \hline
    6 & SQL Queries & 9 \\ \hline
    7 & Project Demonstration & 15 \\ \hline
    8 & Self-learning beyond classroom & 16 \\ \hline
    9 & Learning from the project & 17 \\ \hline
    10 & Challenges faced & 18 \\ \hline
    11 & Conclusion & 19 \\ \hline
\end{tabularx}

\newpage

% --- SECTIONS ---

\section*{I. Storyline}
In the modern industrial landscape, environmental sustainability has become as critical as operational efficiency. Companies face increasing pressure to track and reduce their carbon footprint, especially within complex supply chains. However, most organizations struggle with fragmented data regarding material sources, transport methods, and manufacturing compositions.

\textbf{GreenChain} is a relational database-driven system designed to solve this problem. The story begins with a manufacturing firm that needs to monitor the environmental impact of its products from "cradle to gate."

\textbf{The Workflow:}
\begin{enumerate}
    \item \textbf{Sourcing:} Raw materials (like Steel, Aluminum, or Plastic) are sourced from various \textbf{Suppliers} located in different cities.
    \item \textbf{Transport:} These materials are transported via different \textbf{Modes} (Truck, Ship, Air, Rail). Each shipment is recorded in a \textbf{Transport Log}, capturing the weight and the distance covered.
    \item \textbf{Production:} These materials are then used in specific quantities to manufacture \textbf{Products}. This relationship is managed through a \textbf{Product Composition} registry.
    \item \textbf{Quantification:} By linking these entities, the system calculates the total carbon emission of a product by considering the \textbf{Emission Factor} of each material and the environmental cost of its transportation.
\end{enumerate}

This framework empowers managers to make data-driven decisions, such as choosing suppliers with lower transport distances or switching to materials with lower emission factors.

\newpage

\section*{II. Components of Database Design}

\subsection*{1. Entities and Attributes}

\begin{enumerate}
    \item \textbf{USERS}: Stores authentication and profile details.
    \begin{itemize}
        \item \textbf{Attributes:} \underline{user\_id} (PK), username, password, email, full\_name, created\_at.
    \end{itemize}
    
    \item \textbf{SUPPLIERS}: Represents entities providing raw materials.
    \begin{itemize}
        \item \textbf{Attributes:} \underline{Supplier\_ID} (PK), Supplier\_Name, City, Distance\_KM.
    \end{itemize}
    
    \item \textbf{MATERIALS}: Represents raw materials with their baseline emission impact.
    \begin{itemize}
        \item \textbf{Attributes:} \underline{Material\_ID} (PK), Material\_Name, Emission\_Factor.
    \end{itemize}
    
    \item \textbf{PRODUCTS}: Final goods manufactured by the company.
    \begin{itemize}
        \item \textbf{Attributes:} \underline{Product\_ID} (PK), Product\_Name, Batch\_Number.
    \end{itemize}
    
    \item \textbf{TRANSPORT\_LOG}: Tracks the movement of materials from suppliers.
    \begin{itemize}
        \item \textbf{Attributes:} \underline{Log\_ID} (PK), Material\_ID (FK), Supplier\_ID (FK), Weight\_KG, Transport\_Mode.
    \end{itemize}
    
    \item \textbf{PRODUCT\_COMPOSITION}: Maps materials to products.
    \begin{itemize}
        \item \textbf{Attributes:} \underline{Composition\_ID} (PK), Product\_ID (FK), Material\_ID (FK), Quantity\_Used.
    \end{itemize}
\end{enumerate}

\subsection*{2. Relationships}

\begin{itemize}
    \item \textbf{Supplier to Material (via Transport Log):} Many-to-Many. A supplier can provide many materials, and a material can be sourced from many suppliers.
    \item \textbf{Product to Material (via Product Composition):} Many-to-Many. A product is made of many materials, and a material can be used in many products.
    \item \textbf{Cardinality:}
    \begin{itemize}
        \item One Supplier can have many Transport Logs (1:N).
        \item One Material can appear in many Transport Logs (1:N).
        \item One Product can have many Composition entries (1:N).
    \end{itemize}
\end{itemize}

\newpage

\section*{III. Entity Relationship Diagram}

The ER Diagram for GreenChain serves as the blueprint for our data architecture, visually representing the organizational logic of our sustainability tracking system.

\subsection*{1. Detailed Entity Descriptions}
\begin{itemize}
    \item \textbf{Strong Entities:} 
    \begin{itemize}
        \item \textbf{SUPPLIERS:} Independent entities that exist regardless of other data. They represent the origin point of the supply chain.
        \item \textbf{MATERIALS:} Core building blocks of the products, defined by their chemical or physical properties and static emission factors.
        \item \textbf{PRODUCTS:} The output of the manufacturing process, uniquely identified by batch numbers.
        \item \textbf{USERS:} Manages system access and administrative oversight.
    \end{itemize}
    \item \textbf{Associative Entities (Weak/Relationship-based):}
    \begin{itemize}
        \item \textbf{TRANSPORT\_LOG:} A critical link between Suppliers and Materials. It doesn't just represent a relationship; it captures transaction-specific data like \textit{Weight} and \textit{Transport Mode}.
        \item \textbf{PRODUCT\_COMPOSITION:} Defines the internal recipe of a product, linking specific Materials to Products with a quantitative attribute (\textit{Quantity Used}).
    \end{itemize}
\end{itemize}

\subsection*{2. Relationship Cardinality and Logic}
\begin{itemize}
    \item \textbf{Sourcing Relationship:} One Supplier can have many Transport Logs, but each log entry belongs to exactly one supplier (1:N). This ensures accountability for every kilogram of material moved.
    \item \textbf{Material Utility:} One Material can be used across multiple Products and can be sourced from various Suppliers (M:N). This is resolved through the two associative tables mentioned above.
    \item \textbf{Product Specification:} A single Product batch consists of multiple Material entries (1:N from Product to Composition), allowing for granular carbon footprint calculation at the batch level.
\end{itemize}

\begin{figure}[H]
    \centering
    \includegraphics[width=0.85\textwidth]{er_diagram.png}
    \caption{Entity Relationship Diagram for GreenChain System}
    \label{fig:er_diagram}
\end{figure}

\newpage

\section*{IV. Relational Model}

The transition from the conceptual ER diagram to the logical Relational Model involves strictly following mapping rules to ensure data integrity and query efficiency.

\subsection*{1. Mapping Rules Applied}
\begin{enumerate}
    \item \textbf{Strong Entity Mapping:} Each strong entity from the ER diagram (Users, Suppliers, Materials, Products) was converted into a dedicated relation with its primary key as the unique identifier.
    \item \textbf{Many-to-Many Resolution:} The many-to-many relationship between Products and Materials was converted into an associative table (\texttt{PRODUCT\_COMPOSITION}). This table includes its own surrogate primary key (\texttt{Composition\_ID}) and foreign keys pointing to both participating entities.
    \item \textbf{Transactional Mapping:} The sourcing relationship, which includes attributes like transport mode and weight, was mapped to the \texttt{TRANSPORT\_LOG} table. This allows us to track the "Carbon Intensity of Transit" independently of the material's inherent footprint.
\end{enumerate}

\subsection*{2. Table Schema Definitions}

\begin{enumerate}
    \item \textbf{USERS} (\underline{user\_id}, username, password, email, full\_name, created\_at)
    \item \textbf{SUPPLIERS} (\underline{Supplier\_ID}, Supplier\_Name, City, Distance\_KM)
    \item \textbf{MATERIALS} (\underline{Material\_ID}, Material\_Name, Emission\_Factor)
    \item \textbf{PRODUCTS} (\underline{Product\_ID}, Product\_Name, Batch\_Number)
    \item \textbf{TRANSPORT\_LOG} (\underline{Log\_ID}, \textit{Material\_ID}, \textit{Supplier\_ID}, Weight\_KG, Transport\_Mode)
    \item \textbf{PRODUCT\_COMPOSITION} (\underline{Composition\_ID}, \textit{Product\_ID}, \textit{Material\_ID}, Quantity\_Used)
\end{enumerate}

\subsection*{3. Referential Integrity and Constraints}
\begin{itemize}
    \item \textbf{Primary Keys (PK):} Ensured using \texttt{AUTO\_INCREMENT} for integer IDs to prevent collision and simplify indexing.
    \item \textbf{Foreign Keys (FK):} 
    \begin{itemize}
        \item \texttt{TRANSPORT\_LOG.Material\_ID} $\rightarrow$ \texttt{MATERIALS.Material\_ID}
        \item \texttt{TRANSPORT\_LOG.Supplier\_ID} $\rightarrow$ \texttt{SUPPLIERS.Supplier\_ID}
        \item \texttt{PRODUCT\_COMPOSITION.Product\_ID} $\rightarrow$ \texttt{PRODUCTS.Product\_ID}
        \item \texttt{PRODUCT\_COMPOSITION.Material\_ID} $\rightarrow$ \texttt{MATERIALS.Material\_ID}
    \end{itemize}
    \item \textbf{Delete Cascades:} We implemented \texttt{ON DELETE CASCADE} for the relationship tables. This means if a Product is discontinued (deleted), its composition records are automatically purged to prevent "orphan" data.
\end{itemize}

\newpage

\section*{V. Normalization}

The database design adheres to the following Normal Forms:

\subsection*{1. First Normal Form (1NF)}
All tables contain only atomic values. There are no repeating groups or multi-valued attributes. Each record is unique through a defined Primary Key (e.g., Supplier\_ID, Log\_ID).

\subsection*{2. Second Normal Form (2NF)}
The database is in 1NF and all non-key attributes are fully functionally dependent on the primary key. In the associative tables (\texttt{TRANSPORT\_LOG}, \texttt{PRODUCT\_COMPOSITION}), every attribute depends on the surrogate primary key, and by extension, the composite logic they represent.

\subsection*{3. Third Normal Form (3NF)}
The database is in 2NF and there are no transitive dependencies. For example, in the \texttt{SUPPLIERS} table, \texttt{City} and \texttt{Distance\_KM} depend directly on \texttt{Supplier\_ID}. In \texttt{MATERIALS}, \texttt{Emission\_Factor} depends only on \texttt{Material\_ID}.

\subsection*{Conclusion on Normalization}
The design is highly normalized to 3NF, ensuring data integrity and minimizing redundancy. Since all dependencies are on the key and nothing but the key, it also satisfies BCNF.

\newpage

\section*{VI. SQL Queries}

The system is implemented using \textbf{MySQL 8.0}. Below are the queries used for data definition, manipulation, and analysis.

\subsection*{1. Table Creation}
\begin{lstlisting}
CREATE TABLE MATERIALS (
    Material_ID INT PRIMARY KEY AUTO_INCREMENT,
    Material_Name VARCHAR(255) NOT NULL,
    Emission_Factor DECIMAL(10,2) NOT NULL
);

CREATE TABLE TRANSPORT_LOG (
    Log_ID INT PRIMARY KEY AUTO_INCREMENT,
    Material_ID INT NOT NULL,
    Supplier_ID INT NOT NULL,
    Weight_KG DECIMAL(10,2) NOT NULL,
    Transport_Mode ENUM('Air', 'Ship', 'Truck', 'Rail') NOT NULL,
    FOREIGN KEY (Material_ID) REFERENCES MATERIALS(Material_ID),
    FOREIGN KEY (Supplier_ID) REFERENCES SUPPLIERS(Supplier_ID)
);
\end{lstlisting}

\subsection*{2. Basic Selection}
\textbf{Question:} Retrieve all materials registered in the system.
\begin{lstlisting}
SELECT * FROM MATERIALS;
\end{lstlisting}

\subsection*{3. Filtered Selection}
\textbf{Question:} Find all suppliers located in 'Mumbai'.
\begin{lstlisting}
SELECT Supplier_Name, Distance_KM FROM SUPPLIERS WHERE City = 'Mumbai';
\end{lstlisting}

\subsection*{4. Sorting Results}
\textbf{Question:} List materials in descending order of their emission factors.
\begin{lstlisting}
SELECT Material_Name, Emission_Factor FROM MATERIALS ORDER BY Emission_Factor DESC;
\end{lstlisting}

\subsection*{5. Aggregation: Count}
\textbf{Question:} How many products are currently being tracked?
\begin{lstlisting}
SELECT COUNT(*) AS Total_Products FROM PRODUCTS;
\end{lstlisting}

\subsection*{6. Aggregation: Sum}
\textbf{Question:} What is the total weight of materials transported via 'Truck'?
\begin{lstlisting}
SELECT SUM(Weight_KG) AS Total_Truck_Weight FROM TRANSPORT_LOG WHERE Transport_Mode = 'Truck';
\end{lstlisting}

\subsection*{7. Aggregation: Average}
\textbf{Question:} What is the average distance of all suppliers?
\begin{lstlisting}
SELECT AVG(Distance_KM) AS Avg_Distance FROM SUPPLIERS;
\end{lstlisting}

\subsection*{8. Group By}
\textbf{Question:} Calculate the number of shipments for each transport mode.
\begin{lstlisting}
SELECT Transport_Mode, COUNT(*) AS Shipment_Count 
FROM TRANSPORT_LOG 
GROUP BY Transport_Mode;
\end{lstlisting}

\subsection*{9. Group By with Having}
\textbf{Question:} List transport modes that have moved more than 500 KG in total.
\begin{lstlisting}
SELECT Transport_Mode, SUM(Weight_KG) AS Total_Weight 
FROM TRANSPORT_LOG 
GROUP BY Transport_Mode 
HAVING Total_Weight > 500;
\end{lstlisting}

\subsection*{10. Inner Join}
\textbf{Question:} Display transport logs along with the names of the materials and suppliers.
\begin{lstlisting}
SELECT t.Log_ID, m.Material_Name, s.Supplier_Name, t.Weight_KG 
FROM TRANSPORT_LOG t
INNER JOIN MATERIALS m ON t.Material_ID = m.Material_ID
INNER JOIN SUPPLIERS s ON t.Supplier_ID = s.Supplier_ID;
\end{lstlisting}

\subsection*{11. Left Join}
\textbf{Question:} List all suppliers and the number of transport logs they are associated with (including those with zero).
\begin{lstlisting}
SELECT s.Supplier_Name, COUNT(t.Log_ID) AS Log_Count
FROM SUPPLIERS s
LEFT JOIN TRANSPORT_LOG t ON s.Supplier_ID = t.Supplier_ID
GROUP BY s.Supplier_ID;
\end{lstlisting}

\subsection*{12. Distinct Values}
\textbf{Question:} Get a unique list of cities where suppliers are located.
\begin{lstlisting}
SELECT DISTINCT City FROM SUPPLIERS;
\end{lstlisting}

\subsection*{13. Subquery in FROM}
\textbf{Question:} Find the maximum shipment weight among the 'Truck' shipments using a subquery.
\begin{lstlisting}
SELECT MAX(Weight_KG) FROM (SELECT Weight_KG FROM TRANSPORT_LOG WHERE Transport_Mode = 'Truck') AS TruckWeights;
\end{lstlisting}

\subsection*{14. Wildcard Search (LIKE)}
\textbf{Question:} Find all products that belong to the 'BATCH-2024' series.
\begin{lstlisting}
SELECT * FROM PRODUCTS WHERE Batch_Number LIKE 'BATCH-2024%';
\end{lstlisting}

\subsection*{15. IN Operator}
\textbf{Question:} Retrieve logs for specific transport modes: 'Air' and 'Ship'.
\begin{lstlisting}
SELECT * FROM TRANSPORT_LOG WHERE Transport_Mode IN ('Air', 'Ship');
\end{lstlisting}

\subsection*{16. BETWEEN Operator}
\textbf{Question:} Find materials with an emission factor between 1.0 and 5.0.
\begin{lstlisting}
SELECT * FROM MATERIALS WHERE Emission_Factor BETWEEN 1.0 AND 5.0;
\end{lstlisting}

\subsection*{17. Complex Join: Product Composition Analysis}
\textbf{Question:} List all products along with the total quantity of materials used in them.
\begin{lstlisting}
SELECT p.Product_Name, SUM(pc.Quantity_Used) AS Total_Material
FROM PRODUCTS p
JOIN PRODUCT_COMPOSITION pc ON p.Product_ID = pc.Product_ID
GROUP BY p.Product_ID;
\end{lstlisting}

\subsection*{18. Calculation: Carbon Footprint per Log}
\textbf{Question:} Calculate the estimated emission for each transport log (Weight * Material Emission Factor).
\begin{lstlisting}
SELECT t.Log_ID, m.Material_Name, (t.Weight_KG * m.Emission_Factor) AS Estimated_Emission
FROM TRANSPORT_LOG t
JOIN MATERIALS m ON t.Material_ID = m.Material_ID;
\end{lstlisting}

\subsection*{19. Update Query}
\textbf{Question:} Update the emission factor for 'Steel' due to new environmental data.
\begin{lstlisting}
UPDATE MATERIALS SET Emission_Factor = 1.95 WHERE Material_Name = 'Steel';
\end{lstlisting}

\subsection*{20. Delete Query}
\textbf{Question:} Remove a transport log from the system.
\begin{lstlisting}
DELETE FROM TRANSPORT_LOG WHERE Log_ID = 5;
\end{lstlisting}

\subsection*{21. Existential Check}
\textbf{Question:} Find suppliers who have at least one transport log recorded.
\begin{lstlisting}
SELECT Supplier_Name FROM SUPPLIERS s
WHERE EXISTS (SELECT 1 FROM TRANSPORT_LOG t WHERE t.Supplier_ID = s.Supplier_ID);
\end{lstlisting}

\vspace{1cm}
\section*{VI. Project Demonstration}

\subsection*{1. Tools and Technologies Used}
\begin{itemize}
    \item \textbf{Database:} MySQL 8.0 (Relational Database Management System)
    \item \textbf{Backend:} Node.js with Express.js framework
    \item \textbf{Frontend:} AngularJS (v1.8.2) for the Single Page Application (SPA) architecture
    \item \textbf{Visualization:} Chart.js (v4.4.1) for generating emission analytics and distribution charts
    \item \textbf{Security:} Bcrypt.js for password hashing and Express-Session for session management
    \item \textbf{Styling:} Custom CSS3 with modern gradients and responsive layouts
\end{itemize}

\subsection*{2. GUI Description}
The application features a comprehensive dashboard and management modules:
\begin{itemize}
    \item \textbf{Authentication:} A secure login and registration system.
    \item \textbf{Dashboard:} Real-time visualization of transport mode distribution and top suppliers using Pie and Bar charts.
    \item \textbf{Management Modules:} Dedicated views for Suppliers, Materials, Products, and Transport Logs, allowing full CRUD (Create, Read, Update, Delete) operations.
    \item \textbf{Composition View:} A specialized interface to define the material makeup of each product.
\end{itemize}

\subsection*{3. Screenshots and Snapshots}

The following screenshots illustrate the GreenChain Dashboard and the backend database implementation, showcasing the real-time data visualization and structural integrity of the system.

\begin{figure}[H]
    \centering
    \begin{minipage}{0.48\textwidth}
        \centering
        \includegraphics[width=\textwidth]{image.png}
        \caption{Dashboard Overview: Displaying total counts for Suppliers, Materials, and Products.}
        \label{fig:dashboard_overview}
    \end{minipage}
    \hfill
    \begin{minipage}{0.48\textwidth}
        \centering
        \includegraphics[width=\textwidth]{image copy.png}
        \caption{Analytics Visualization: Interactive Pie Chart for Transport Mode Distribution.}
        \label{fig:dashboard_analytics}
    \end{minipage}
\end{figure}

\begin{figure}[H]
    \centering
    \includegraphics[width=0.8\textwidth]{image copy 2.png}
    \caption{System Navigation: Accessible modules for comprehensive supply chain management.}
    \label{fig:dashboard_navigation}
\end{figure}

\begin{figure}[H]
    \centering
    \includegraphics[width=0.8\textwidth]{image copy 3.png}
    \caption{Database Schema Verification: Terminal output showing the successfully created tables and their structures within the MySQL environment.}
    \label{fig:database_schema}
\end{figure}

\vspace{1cm}
\section*{VII. Self -Learning beyond classroom}

Throughout the development of GreenChain, several concepts were explored beyond the standard curriculum:

\begin{enumerate}
    \item \textbf{Object-Relational Mapping (ORM) vs. Raw Queries:} While we used the \texttt{mysql2} driver for raw SQL queries to better understand the database logic, we learned about the trade-offs between raw SQL and ORMs like Sequelize.
    \item \textbf{Asynchronous Database Operations:} Implementing database calls in Node.js required a deep dive into Promises and \texttt{async/await} patterns to handle the non-blocking nature of the environment.
    \item \textbf{Security Best Practices:} We learned about SQL Injection prevention by using prepared statements and parameterized queries (e.g., using \texttt{?}) instead of string concatenation.
    \item \textbf{Data Visualization:} Integrating \texttt{Chart.js} taught us how to transform raw SQL aggregate results (JSON) into visual formats that provide business value.
    \item \textbf{State Management:} Using Express sessions to maintain user state across different routes in a web application.
\end{enumerate}

\vspace{1cm}
\section*{VIII. Learning from the Project}

This project has been an invaluable learning experience in bridging the gap between theoretical database concepts and real-world application development:

\begin{itemize}
    \item \textbf{Schema Design:} We realized that the most important part of a project happens before a single line of code is written—the design of the ER diagram. A well-designed schema makes development significantly easier.
    \item \textbf{SQL Proficiency:} Writing over 20 diverse queries helped solidify our understanding of Joins, Subqueries, and Grouping.
    \item \textbf{Problem Solving:} It taught us how to model complex real-world relationships (like the carbon impact of transport) into clean, relational tables.
    \item \textbf{Full-Stack Integration:} We understood how a backend API acts as a bridge between the frontend user interface and the persistent storage (MySQL).
\end{itemize}

\vspace{1cm}
\section*{IX. Challenges Faced}

\begin{enumerate}
    \item \textbf{Database Connection Management:} Initially, we faced issues with "Too many connections." This was resolved by implementing a connection pool using \texttt{mysql.createPool}.
    \item \textbf{Complex Joins:} Crafting a query that could calculate the total emission of a product by joining \texttt{PRODUCTS}, \texttt{COMPOSITION}, \texttt{MATERIALS}, and \texttt{TRANSPORT\_LOG} was challenging but rewarding.
    \item \textbf{Data Consistency:} Ensuring that materials couldn't be deleted if they were already part of a product composition required proper implementation of Foreign Key constraints and \texttt{ON DELETE} strategies.
    \item \textbf{Frontend Integration:} Mapping the nested JSON responses from complex SQL joins to the AngularJS view layer required careful data transformation.
    \item \textbf{Environment Configuration:} Setting up the local MySQL environment across different operating systems (macOS and Windows) required careful handling of root permissions and port configurations.
\end{enumerate}

\vspace{1cm}
\section*{X. Conclusion}

The \textbf{GreenChain} project successfully demonstrates the power of Relational Database Management Systems in solving modern business and environmental challenges. Key takeaways include:

\begin{itemize}
    \item \textbf{Data Integrity:} RDBMS ensures that our supply chain records remain consistent and reliable through constraints and normalization.
    \item \textbf{Scalability:} The normalized design allows the system to grow as more materials and suppliers are added without redundant data storage.
    \item \textbf{Informed Decision Making:} By using SQL aggregations, the system provides actionable insights into which parts of the supply chain are most damaging to the environment.
\end{itemize}

In conclusion, this project highlights how database technology can be leveraged to drive sustainability and operational excellence in the industrial sector.

\end{document}
